% !TeX program = xelatex
% Самодостатній файл UTF-8. Компіляція: xelatex <ім’я-файлу>.tex (двічі).
\documentclass[10pt,a4paper,landscape]{article}
\usepackage[margin=15mm]{geometry}
\usepackage{fontspec}
\setmainfont{FreeSerif.otf}[BoldFont=FreeSerifBold.otf,ItalicFont=FreeSerifItalic.otf,BoldItalicFont=FreeSerifBoldItalic.otf]
\usepackage{polyglossia}
\setdefaultlanguage{ukrainian}
\usepackage{longtable,array,booktabs,xcolor,fancyhdr}
\usepackage[hidelinks]{hyperref}
\definecolor{accent}{HTML}{16656B}
\pagestyle{fancy}\fancyhf{}
\fancyhead[L]{STEM • 6 клас}\fancyhead[R]{Навчальний та календарно-тематичний план}
\fancyfoot[C]{\thepage}
\setlength{\headheight}{14pt}
\setlength{\parindent}{0pt}
\setlength{\parskip}{5pt}
\renewcommand{\arraystretch}{1.18}
\newcolumntype{P}[1]{>{\raggedright\arraybackslash}p{#1}}
\begin{document}
{\LARGE\color{accent} Навчальний план STEM, 6 клас}

{\large 1 год/тиждень • 32 години на рік}

Заклад освіти: \rule{90mm}{0.3pt}\quad Навчальний рік: \rule{30mm}{0.3pt}

Учитель/учителька: \rule{95mm}{0.3pt}\quad Клас: \rule{25mm}{0.3pt}
\section*{Пояснювальна записка}
Авторський адаптований план укладено за наданою модельною програмою «STEM. 5--6 класи (міжгалузевий інтегрований курс)» О. В. Бутурліної та О. Є. Артєм’євої; розділ 6 класу — с. 21--32 PDF «STEM.5-6.kl.Buturlina.Artyemyeva.04.10.pdf». Програма є основою добору змісту й очікуваних результатів. Сканованого зошита для 6 класу не надано, тому номери його сторінок і вправ не зазначаються. Запропоновані практичні роботи є адаптаціями видів діяльності з програми.

Мета: розвивати дослідницькі, інженерні, математичні, цифрові та комунікативні вміння через проєкти, самопізнання й знайомство з професіями. Збережено вступ, п’ять змістових модулів і підсумковий хакатон, фестиваль та STEM-практику.

Орієнтовний розподіл програми: 3 + 7 + 5 + 6 + 5 + 5 + 4 = 35 годин. На с. 8 PDF передбачено самостійний розподіл годин учителем зі збереженням очікуваних результатів. У базовому плані на 32 тижні розподіл становить 3 + 6 + 6 + 5 + 4 + 4 + 4: споріднені теми об’єднано, а для тривалих спостережень за рослинами передбачено шість тижнів. Результатом кожного модуля є завершений проєкт із презентацією та рефлексією.

\textbf{Облік часу:} І семестр — $15\times1=15$ год; ІІ семестр — $17\times1=17$ год; разом 32 години. Один рядок таблиці відповідає одному уроку. Для збереження практичної роботи використовуються прості моделі, готові шаблони та підготовлені вчителем матеріали. Обов’язкові завдання виконуються в навчальний час; обсяг виробу відповідає можливостям одного уроку.

\textbf{Організація спостережень:} дослід із рослинами починається на тижні 10 та завершується на тижні 15. На наступних уроках цього модуля 5--10 хвилин відводиться на журнал спостережень і догляд; цей час входить у вказані години. Між уроками догляд організовує вчитель або добровільна чергова група. У разі невдалого пророщування використовуються готові саджанці чи заздалегідь підготовлені дані; це не замінює аналізу причин і результатів.

\textbf{Ресурси та безпека:} папір, картон, лінійки, циркуль, секундомір, термометр, компас, безпечні рослини, ємності, нитки, готові деталі моделей; за можливості — камера, комп’ютер, датчик вологості та низьковольтний навчальний набір. Цифрові завдання мають паперову альтернативу. Роботи з водою відокремлено від електричних з’єднань; тільки батарейне живлення. Не використовуються токсичні рослини, пестициди, полум’я чи піротехнічні ракети. Запуски моделей — у вільній зоні під наглядом учителя. Для медіаробіт застосовуються власні або дозволені матеріали; публічні акаунти учнів не потрібні.

\textbf{Джерельні позначення:} у таблиці вказані сторінки PDF навчальної програми, а не зошита. Матеріали про підприємства, професії та медіа добирає і перевіряє вчитель перед заняттям. При дослідженні золотого перерізу учні перевіряють припущення вимірюваннями; подібність до спіралі або пропорції не доводить універсальної закономірності.
\newpage
\section*{Розподіл навчального часу}
\begin{tabular}{P{155mm}rrr}
\toprule Розділ & І семестр & ІІ семестр & Разом \\\midrule
Вступ. Україна та мій регіон & 3 & 0 & 3 \\
Людина — людина. Пізнай себе та світ навколо себе & 6 & 0 & 6 \\
Людина — природа. Сад на підвіконні & 6 & 0 & 6 \\
Людина — техніка. Від возу до космічної ракети & 0 & 5 & 5 \\
Людина — образ. Я так бачу! & 0 & 4 & 4 \\
Людина — знак. Послідовність Фібоначчі. Золотий переріз & 0 & 4 & 4 \\
Хакатон, фестиваль, STEM-практика & 0 & 4 & 4 \\
\midrule Разом & 15 & 17 & 32 \\
\bottomrule\end{tabular}
\section*{Оцінювання та очікувані результати}
На першому уроці проводяться короткі стартові завдання на читання карти, таблиці та планування дослідження. Формувальне оцінювання триває під час кожного проєкту: спостереження, записи даних, ескізи, моделі, пояснення, самооцінювання та взаємний відгук. Діагностика за І семестр інтегрована у тиждень 15, річне підсумкове оцінювання — у тиждень 32; окремих додаткових годин не потрібно.

Спільні критерії: зрозуміла проблема й мета; коректні вимірювання та обчислення; обґрунтоване рішення; функціональність і безпечність моделі; аналіз випробувань; дотримання авторства; співпраця та особистий внесок. Формувальний відгук описує виконання з допомогою, самостійне виконання або перенесення вміння у нову ситуацію. Конкретну шкалу підсумкового оцінювання визначає заклад освіти.

Наприкінці року учень/учениця пов’язує професії з ресурсами регіону й власними вміннями; планує командну роботу; досліджує умови росту рослин і пояснює алгоритм автополиву; вимірює шлях та час, обчислює швидкість і створює транспортну модель; створює медіапродукт із дотриманням етики; застосовує правило Фібоначчі та перевіряє пропорції; презентує й удосконалює проєкт на основі даних.

Портфоліо містить карту професій регіону, навчальне резюме, спільний проєкт класу, журнал рослин і модель поливу, транспортну модель, медіароботу, математично-художню композицію та підсумковий прототип. Посильний рівень забезпечують шаблони й приклади; поглиблення — самостійний вибір змінної, повторні вимірювання та аргументоване вдосконалення.
\newpage
\section*{Календарно-тематичний план. І семестр}
\small\begin{longtable}{P{13mm}P{16mm}P{47mm}P{18mm}P{88mm}P{52mm}}
\toprule Тиждень / год & Дата & Модуль і тема & Програма PDF, с. & Навчальна діяльність & Очікуваний результат \\\midrule\endfirsthead
\toprule Тиждень / год & Дата & Модуль і тема & Програма PDF, с. & Навчальна діяльність & Очікуваний результат \\\midrule\endhead
\bottomrule\endfoot
1 / 1 & & \textbf{Вступ}\par Україна у світовій економіці. Людський капітал & 21--22 & Стартові завдання на читання даних; карта STEM-галузей; обговорення ролі знань і праці. & Пояснює поняття людського капіталу та пов’язує галузі з професіями. \\ \midrule
2 / 1 & & \textbf{Вступ}\par Мій регіон: ресурси, підприємства та професії & 21--22 & Позначити на карті адміністративний центр і підприємства; визначити координати та відстань за масштабом. & Пов’язує ресурси регіону з видами діяльності; читає карту й обчислює відстань. \\ \midrule
3 / 1 & & \textbf{Вступ}\par Професії родини. Династії та моє майбутнє & 21--22 & Створити дерево професій за добровільно наданими або вигаданими даними; визначити власну навчальну ціль. & Пояснює поняття професійної династії та зв’язок освіти з майбутньою діяльністю. \\ \midrule
4 / 1 & & \textbf{Людина — людина}\par Я у світі людей: таланти та здібності & 22--23 & Скласти карту інтересів, досягнень і сильних сторін; обговорити приклади розвитку здібностей. & Називає власні сильні сторони та підкріплює їх прикладами. \\ \midrule
5 / 1 & & \textbf{Людина — людина}\par Навички XXI століття. Hard skills і soft skills & 22--23 & Розподілити картки вмінь; зіставити вимоги професій майбутнього з навчальними діями. & Розрізняє фахові й соціальні навички та пояснює їх поєднання. \\ \midrule
6 / 1 & & \textbf{Людина — людина}\par Моє перше резюме та портфоліо & 22--23 & Заповнити навчальне резюме за шаблоном; дібрати приклади робіт до паперового або цифрового портфоліо. & Структурує відомості про вміння й досягнення без зайвих персональних даних. \\ \midrule
7 / 1 & & \textbf{Людина — людина}\par Професія психолога. Запитання для інтерв’ю & 22--23 & Підготувати запитання до шкільного психолога; провести зустріч або опрацювати професійний кейс. & Описує завдання, освіту й навички психолога; ставить доречні запитання. \\ \midrule
8 / 1 & & \textbf{Людина — людина}\par Команда, ролі та лідерство & 22--23 & Виконати командне конструкторське завдання; змінити ролі; обговорити способи досягнення згоди. & Розподіляє роботу, слухає партнерів і оцінює власний внесок. \\ \midrule
9 / 1 & & \textbf{Людина — людина}\par Проєкт «Галактика 6-го класу» & 22--23 & Об’єднати інтереси та вміння класу в спільне портфоліо «Ми різні — ми рівні»; представити результат. & Створює спільний продукт і пояснює цінність різних здібностей у команді. \\ \midrule
10 / 1 & & \textbf{Людина — природа}\par Кімнатні рослини. Плануємо сад на підвіконні & 24--25 & За атласом дібрати безпечні рослини; закласти дослід із насінням або готовими саджанцями; розпочати журнал. & Розпізнає вибрані рослини й формулює умови дослідження. \\ \midrule
11 / 1 & & \textbf{Людина — природа}\par Умови росту: світло, температура, вода & 24--25 & Визначити орієнтацію вікна компасом; виміряти температуру; порівняти умови, змінюючи один чинник. & Добирає спосіб вимірювання та фіксує умови й результати. \\ \midrule
12 / 1 & & \textbf{Людина — природа}\par Догляд за рослинами. Простий флораріум & 24--25 & Висадити вибрані рослини або створити відкритий флораріум; скласти інструкцію догляду. & Дотримується послідовності висаджування та пояснює потреби рослини. \\ \midrule
13 / 1 & & \textbf{Людина — природа}\par Автополив: задум, схема, матеріали & 24--25 & Побудувати схему подачі води; зібрати просту модель ґнотового або крапельного поливу. & Пояснює принцип роботи моделі та добирає матеріали за призначенням. \\ \midrule
14 / 1 & & \textbf{Людина — природа}\par Розумна теплиця. Датчики й алгоритм керування & 24--25 & На демонстрації простежити роботу датчика, контролера й помпи; скласти алгоритм «якщо — то». & Розрізняє вимірювання, прийняття рішення та виконання команди. \\ \midrule
15 / 1 & & \textbf{Людина — природа}\par Сад на підвіконні: захист проєкту & 24--25 & Порівняти записи спостережень; презентувати систему догляду й картку професії; провести семестрову діагностику. & Підтверджує висновок даними; пов’язує проєкт із ботанікою, квітникарством і фітодизайном. \\ \midrule
\end{longtable}\normalsize
\newpage
\section*{Календарно-тематичний план. ІІ семестр}
\small\begin{longtable}{P{13mm}P{16mm}P{47mm}P{18mm}P{88mm}P{52mm}}
\toprule Тиждень / год & Дата & Модуль і тема & Програма PDF, с. & Навчальна діяльність & Очікуваний результат \\\midrule\endfirsthead
\toprule Тиждень / год & Дата & Модуль і тема & Програма PDF, с. & Навчальна діяльність & Очікуваний результат \\\midrule\endhead
\bottomrule\endfoot
16 / 1 & & \textbf{Людина — техніка}\par Колесо в історії техніки. Радіус і довжина кола & 26--27 & Побудувати коло циркулем і ниткою; виміряти радіус та довжину кола; скласти колісну модель. & Знаходить зв’язок між обертом колеса й пройденим шляхом. \\ \midrule
17 / 1 & & \textbf{Людина — техніка}\par Рух: траєкторія, шлях, переміщення, швидкість & 26--27 & Дослідити рух моделі на розміченій доріжці; виміряти шлях і час; порівняти рівномірний і нерівномірний рух. & Розрізняє шлях і переміщення; обчислює середню швидкість; описує зміну швидкості. \\ \midrule
18 / 1 & & \textbf{Людина — техніка}\par Енергія транспорту. Реактивний рух & 26--27 & Обговорити рух тіла й передавання тепла; випробувати кульку на нитці; пояснити реактивний рух і космічну швидкість якісно. & Пов’язує джерело енергії з рухом та пояснює принцип реактивної моделі. \\ \midrule
19 / 1 & & \textbf{Людина — техніка}\par Транспортні системи та інфраструктура & 26--27 & Порівняти види транспорту; побудувати маршрут; дібрати професії й запропонувати екологічне поліпшення. & Обґрунтовує вибір транспорту з огляду на відстань, потреби та ресурси. \\ \midrule
20 / 1 & & \textbf{Людина — техніка}\par Проєкт «Транспорт майбутнього» & 26--27 & За ескізом створити модель із готових деталей; випробувати рух і представити задум альтернативного приводу. & Демонструє модель, фіксує результати й аргументує вдосконалення. \\ \midrule
21 / 1 & & \textbf{Людина — образ}\par Статичні образи: від малюнка до цифрового зображення & 27--29 & Створити образ різними засобами; порівняти фотографію, малюнок та приклад голограми; обговорити оптичну ілюзію. & Розрізняє способи фіксації зображення та добирає засіб до задуму. \\ \midrule
22 / 1 & & \textbf{Людина — образ}\par Кіно, анімація, комп’ютерна графіка & 27--29 & Створити коротку покадрову або GIF-анімацію; за відсутності техніки — фліпбук; порівняти реальні й віртуальні образи. & Пояснює послідовність кадрів і створює простий рухомий образ. \\ \midrule
23 / 1 & & \textbf{Людина — образ}\par Мова медіа. Соціальні мережі та відповідальність & 27--29 & Проаналізувати повідомлення; перевірити авторство; скласти правила приватності, мережевого етикету та захисту репутації. & Відрізняє факт від оцінки; добирає безпечний спосіб поширення медіа. \\ \midrule
24 / 1 & & \textbf{Людина — образ}\par Проєкт «Я так бачу!» & 27--29 & Оформити фотосерію, анімацію або графічну історію «Погляд через об’єктив»; представити задум і пов’язані професії. & Поєднує виражальні засоби з повідомленням і зазначає авторів використаних матеріалів. \\ \midrule
25 / 1 & & \textbf{Людина — знак}\par Послідовність Фібоначчі. Правило й алгоритм & 29--30 & Побудувати послідовність від 0 і 1; обчислити наступні члени; записати алгоритм словами або в середовищі програмування. & Застосовує правило суми двох попередніх членів і знаходить член за номером. \\ \midrule
26 / 1 & & \textbf{Людина — знак}\par Квадрати, спіралі та закономірності у природі & 29--30 & Побудувати квадрати зі сторонами за числами Фібоначчі; накреслити дуги; дослідити підготовлені зображення рослин. & Пов’язує числовий ряд із геометричною моделлю; відрізняє модель від природного об’єкта. \\ \midrule
27 / 1 & & \textbf{Людина — знак}\par Золотий переріз. Перевіряємо пропорції & 29--30 & Порівняти відношення сусідніх чисел; виміряти прямокутники; перевірити твердження про пропорції будівель і творів мистецтва. & Обчислює відношення, порівнює його з наближенням 1,6 та оцінює точність висновку. \\ \midrule
28 / 1 & & \textbf{Людина — знак}\par Проєкт «Пропорції у природі та естетиці» & 29--30 & Створити логотип або композицію з поясненням пропорцій; додати ритмічний чи поетичний фрагмент за числами Фібоначчі. & Застосовує математичне правило у творчій роботі та пояснює межі обраної моделі. \\ \midrule
29 / 1 & & \textbf{Підсумкова STEM-практика}\par STEM-практика: професії та підприємства громади & 31--32 & Провести очну або віртуальну зустріч із фахівцем; за відсутності — дослідити підготовлений кейс. & Ставить змістовні запитання й пов’язує професію з навчальними вміннями. \\ \midrule
30 / 1 & & \textbf{Підсумкова STEM-практика}\par Хакатон: проблема, план і прототип & 31--32 & Обрати проблему класу або громади; визначити критерії успіху, ролі й ресурси; зібрати простий прототип. & Планує послідовність дій і добирає дані для розв’язання проблеми. \\ \midrule
31 / 1 & & \textbf{Підсумкова STEM-практика}\par Хакатон: випробування та вдосконалення & 31--32 & Випробувати прототип за узгодженими критеріями; записати результати; підготувати короткий захист. & Обґрунтовує зміни даними, враховує відгуки й оцінює ресурси. \\ \midrule
32 / 1 & & \textbf{Підсумкова STEM-практика}\par STEM-фестиваль. Підсумкова рефлексія & 31--32 & Представити портфоліо й прототип; виконати підсумкові завдання та самооцінювання. & Демонструє поступ у дослідженні, конструюванні, роботі з медіа та співпраці. \\ \midrule
\end{longtable}\normalsize
\end{document}
